\begin{frame}{讲解：这门课真正分析的对象是什么}
  \Core{《大模型推理引擎入门》不是一门“模型结构课”，而是一门“在线推理系统课”。}
  \begin{itemize}
    \item 我们关心的是请求怎样进入系统、怎样被调度、怎样携带状态、怎样形成用户看到的时延与吞吐。
    \item 同一个模型参数，在不同 runtime、不同调度器、不同 KV 管理策略下，会表现出完全不同的服务品质。
    \item 因此课程的基本分析单位不是某一层网络，而是“请求对象 + 控制点 + 状态对象 + 结果指标”。
  \end{itemize}
\end{frame}

\begin{frame}{讲解：导论课应带走的统一判断框架}
  \Core{从第一讲开始就要养成同一套判断顺序，否则后面会把现象、机制和结论混在一起。}
  \begin{enumerate}
    \item 先问：系统里长期存在的对象是什么，例如请求、KV、队列、会话状态。
    \item 再问：谁真正改变了系统行为，例如 scheduler、block manager、executor。
    \item 再问：收益如何传递到端到端指标，例如 TTFT、P99、goodput。
    \item 最后问：哪些日志、实验和代码证据能支持这条因果链。
  \end{enumerate}
\end{frame}

\begin{frame}{例子：为什么传统部署经验在 LLM 时代失效（1）}
  \Core{从分类模型到生成模型，服务系统的主矛盾从“单次前向吞吐”转向“跨轮状态组织”。}
  \begin{block}{阶段一：传统 CV / NLP 服务}
    单请求通常是一次前向，输入输出短，显存主要被模型参数占据。
  \end{block}
\end{frame}

\begin{frame}{例子：为什么传统部署经验在 LLM 时代失效（2）}
  \Core{从分类模型到生成模型，服务系统的主矛盾从“单次前向吞吐”转向“跨轮状态组织”。}
  \begin{block}{阶段二：早期小模型对话服务}
    即便开始有生成过程，请求状态较小，系统还能把问题近似看成短任务排队。
  \end{block}
\end{frame}

\begin{frame}{例子：为什么传统部署经验在 LLM 时代失效（3）}
  \Core{从分类模型到生成模型，服务系统的主矛盾从“单次前向吞吐”转向“跨轮状态组织”。}
  \begin{block}{阶段三：大模型在线服务}
    长上下文、连续 decode、KV 驻留和动态到达同时出现，服务问题不再是简单部署问题。
  \end{block}
\end{frame}

\begin{frame}{例子：真实服务里最先撞上的常见问题（1）}
  \Core{很多系统失效并不是“模型不会生成”，而是系统无法稳定承载状态和并发。}
  \begin{block}{现象一：晚高峰突然超时}
    请求并发上升后，TTFT 和 P99 同时恶化，用户感知是“模型卡住了”。
  \end{block}
\end{frame}

\begin{frame}{例子：真实服务里最先撞上的常见问题（2）}
  \Core{很多系统失效并不是“模型不会生成”，而是系统无法稳定承载状态和并发。}
  \begin{block}{现象二：显存没有满却接不了更多请求}
    连续空间、碎片、预留和块管理会让系统在理论容量之前就先失去弹性。
  \end{block}
\end{frame}

\begin{frame}{例子：真实服务里最先撞上的常见问题（3）}
  \Core{很多系统失效并不是“模型不会生成”，而是系统无法稳定承载状态和并发。}
  \begin{block}{现象三：换一套 runtime 行为完全不同}
    同一模型在不同服务栈上的表现差异，说明问题核心在系统组织而不在模型参数本身。
  \end{block}
\end{frame}

\begin{frame}{代码例子：最小请求从哪里进入系统（1）}
  \Core{课程不从复杂优化开始，而是先从一条最短代码路径建立控制点意识。}
  \begin{block}{入口代码}
    {\ttfamily\small
    \detokenize{from nanovllm import LLM, SamplingParams}\par
    \detokenize{llm = LLM(model_path, enforce_eager=True)}\par
    \detokenize{outputs = llm.generate(prompts, sampling_params)}\par
    }
  \end{block}
\end{frame}

\begin{frame}{代码例子：最小请求从哪里进入系统（2）}
  \Core{课程不从复杂优化开始，而是先从一条最短代码路径建立控制点意识。}
  \begin{block}{课堂应追问}
    这里的 \texttt{LLM.generate} 是最终控制点，还是对更深层 engine 的封装入口？
  \end{block}
\end{frame}

\begin{frame}{代码例子：最小请求从哪里进入系统（3）}
  \Core{课程不从复杂优化开始，而是先从一条最短代码路径建立控制点意识。}
  \begin{block}{系统含义}
    只要这一步还说不清，后续调度、KV、实验设计都会失去锚点。
  \end{block}
\end{frame}

\begin{frame}{代码例子：最小日志为什么足够有用（1）}
  \Core{导论课最有价值的代码例子不是复杂优化，而是能证明主路径存在的最小插桩。}
  \begin{block}{最小日志样式}
    {\ttfamily\small
    \detokenize{print(request_id, prompt_len, stage, timestamp)}\par
    \detokenize{print(batch_size, running_size, finished_count)}\par
    }
  \end{block}
\end{frame}

\begin{frame}{代码例子：最小日志为什么足够有用（2）}
  \Core{导论课最有价值的代码例子不是复杂优化，而是能证明主路径存在的最小插桩。}
  \begin{block}{为什么不一开始就打满日志}
    如果日志过多，学生会看到很多现象，但抓不住“谁拥有控制权”这个主问题。
  \end{block}
\end{frame}

\begin{frame}{代码例子：最小日志为什么足够有用（3）}
  \Core{导论课最有价值的代码例子不是复杂优化，而是能证明主路径存在的最小插桩。}
  \begin{block}{讲授重点}
    先用少量证据锁定路径，再逐步增加探针，这是系统调试和研究的共同方法。
  \end{block}
\end{frame}

\begin{frame}{反例：把系统问题误写成模型问题会怎样（1）}
  \Core{导论课最需要纠正的，就是把一切现象都归因于模型大小或算子快慢。}
  \begin{block}{误读一：输出慢是模型不够强}
    实际上，很多慢来自排队、prefill 压力和状态驻留，而不是模型回答质量。
  \end{block}
\end{frame}

\begin{frame}{反例：把系统问题误写成模型问题会怎样（2）}
  \Core{导论课最需要纠正的，就是把一切现象都归因于模型大小或算子快慢。}
  \begin{block}{误读二：换更快 kernel 就能解决服务问题}
    如果请求组织和状态管理失衡，再快的单点算子也难兑现到端到端收益。
  \end{block}
\end{frame}

\begin{frame}{反例：把系统问题误写成模型问题会怎样（3）}
  \Core{导论课最需要纠正的，就是把一切现象都归因于模型大小或算子快慢。}
  \begin{block}{导论阶段的任务}
    先把问题放回系统边界中，再谈优化对象和收益链路。
  \end{block}
\end{frame}

\begin{frame}{课堂练习：把抽象课程问题落到具体对象（1）}
  \Core{学生需要尽早学会把“大问题”翻译成“对象、控制点、证据”三件事。}
  \begin{block}{练习一}
    给出一个你见过的“模型卡顿”现象，判断它更像是排队、状态还是执行路径问题。
  \end{block}
\end{frame}

\begin{frame}{课堂练习：把抽象课程问题落到具体对象（2）}
  \Core{学生需要尽早学会把“大问题”翻译成“对象、控制点、证据”三件事。}
  \begin{block}{练习二}
    在 \texttt{example.py} 到 \texttt{llm\_engine.py} 之间，指出最值得先读的两个函数。
  \end{block}
\end{frame}

\begin{frame}{课堂练习：把抽象课程问题落到具体对象（3）}
  \Core{学生需要尽早学会把“大问题”翻译成“对象、控制点、证据”三件事。}
  \begin{block}{练习三}
    用一句话回答：为什么这门课要从“系统边界”而不是从“某篇论文机制”讲起？
  \end{block}
\end{frame}
